\section{Журнал выполнения}
\begin{enumerate}
    \item Стемминг нужен и в индексе, и в запросе, иначе основы в словаре не совпадут со словами запроса. Чтобы поиск не надо было запускать с отдельным ключом (и нельзя было ошибиться, применив стемминг не к тому индексу), индекс сам сообщает о себе флагом в заголовке. Разбор запроса и операторы ЛР8 не менялись: стеммируется только каждый терм после токенизации. Автотесты ЛР7 и ЛР8 после изменений пройдены.
    \item Стеммер Портера рассчитан на английские слова из букв \texttt{a--z}. В корпусе есть числа, слова с диакритикой и не латиница (ЛР3), поэтому они не стеммируются, иначе правила для гласных и согласных давали бы бессмыслицу.
    \item Реализован алгоритм по статье 1980 года (у более поздних версий Портера несколько правил отличаются, например \emph{-abli} / \emph{-bli}). Проверка --- 77 слов с известными основами из статьи и стандартного словаря примеров (тест~1).
    \item Индексация со стеммингом медленнее примерно в 1,4 раза (таблица ниже): стеммер вызывается для каждого из 36,7 млн токенов, хотя различных токенов всего 775 тыс. Быстрее было бы стеммировать каждое слово один раз --- держать в словаре и токен, и его основу.
    \item Оценка: пул --- 225 документов. По фрагментам не всегда видно, о чём документ, поэтому при сомнении читался полный текст. Общие правила разметки: упоминание вскользь или ссылка на другое место (\enquote{Go next}) --- 0; место из запроса обязательно (рифы Индии для \texttt{snorkeling coral reefs thailand} --- 0). По запросу \texttt{travel news} ни один документ не содержит собственно новостей о путешествиях; релевантными сочтены путеводители, где перечислены местные новостные издания (газеты, радио) для приезжих.
\end{enumerate}

\section{План тестирования}
Тесты --- юнит-тесты на C++ (\texttt{tests/test\_stemmer.cpp}) на собственном минимальном каркасе \texttt{lab3/tests/unittest.h} (макросы \texttt{TEST} и \texttt{CHECK}, без сторонних библиотек). Они вызывают функции программы напрямую, а последний тест запускает утилиты командной строки как отдельные процессы. Сборка и запуск --- \texttt{make test}. Эталонные значения в тестах считаются независимо от проверяемого кода, для этого используются контейнеры STL (\texttt{std::map}, \texttt{std::set}) --- это код тестов, а не поисковой системы.

\begin{center}
\small
\begin{tabular}{|c|p{0.3\textwidth}|p{0.58\textwidth}|}
\hline
\textnumero & Тест & Что проверяется \\
\hline
1 & Шаг 1a & Множественное число: \emph{caresses} $\to$ \emph{caress}, \emph{ponies} $\to$ \emph{poni}, \emph{cats} $\to$ \emph{cat}. \\
\hline
2 & Шаг 1b & \emph{-eed}, \emph{-ed}, \emph{-ing} и восстановление основы: \emph{agreed} $\to$ \emph{agre}, \emph{hopping} $\to$ \emph{hop}, \emph{hiking} $\to$ \emph{hike}, \emph{sing} без изменений. \\
\hline
3 & Шаг 1c & \emph{happy} $\to$ \emph{happi}, \emph{sky} без изменений (нет гласной в основе). \\
\hline
4 & Шаги 2 и 3 & Двойные суффиксы при $m > 0$: \emph{relational} $\to$ \emph{relat}, \emph{generalization} $\to$ \emph{gener}, \emph{goodness} $\to$ \emph{good}. \\
\hline
5 & Шаги 4 и 5 & Суффиксы при $m > 1$, \emph{-ion} только после \emph{s}/\emph{t}, конечная \emph{e} и \emph{ll}: \emph{adoption} $\to$ \emph{adopt}, \emph{rate}, \emph{controlling} $\to$ \emph{control}. \\
\hline
6 & Эталонный словарь & Все 77 слов из \texttt{tests/porter\_cases.tsv} дают известные основы. \\
\hline
7 & Что не меняется & Кириллица, токены с цифрами и диакритикой, слова из 1--2 букв остаются как есть. \\
\hline
8 & Индекс со стеммингом & Сборщик индекса ЛР7 с флагом стемминга на 5 документах: \emph{museum/museums}, \emph{skiing/skis}, \emph{news/new} дают по одному терму с объединённым списком; флаг в заголовке 1, у обычного индекса --- 0. \\
\hline
9 & Поиск & Разбор запроса ЛР8 по индексу со стеммингом приводит слово к основе и находит все формы, операторы работают; по обычному индексу --- только точная форма. \\
\hline
10 & Утилиты & \texttt{build/stem}, \texttt{indexer -s} и \texttt{search} как отдельные процессы. \\
\hline
\end{tabular}
\end{center}

Скрипт оценки (\texttt{scripts/evaluate.py}) завершается с ошибкой, если в первых 10 результатах какой-либо системы есть документ без оценки, --- так P@10 всегда считается по полной разметке.

\section{Результаты}
\subsection{Автотесты}
Пройдены все 10 тестов. Журнал прогона (\texttt{results/tests.log}):

{\scriptsize
\VerbatimInput[breaklines, breakanywhere]{../results/tests.log}
}

\subsection{Индекс}
\begin{center}
\small
\begin{tabular}{|l|r|r|r|}
\hline
Показатель & Без стемминга & Со стеммингом & Изменение \\
\hline
Термов & 775\,205 & 713\,920 & -7,9\,\% \\
\hline
Записей в списках документов & 13\,644\,680 & 12\,562\,895 & -7,9\,\% \\
\hline
Средняя длина терма, символов & 7,87 & 7,67 & -2,5\,\% \\
\hline
Размер индекса, МиБ & 74,0 & 68,2 & -7,8\,\% \\
\hline
Время индексации, с & 7,48 & 10,76 & 43,8\,\% \\
\hline
\end{tabular}

\end{center}
Словарь сократился всего на 8\,\%: большинство термов корпуса --- редкие имена собственные, числа и иностранные слова, у которых нет других форм. В корпусе 39\,095 основ объединяют две и более формы, всего 100\,380 из 775\,205 термов (\texttt{results/conflation.txt}). Какие формы склеиваются для основ из запросов:
\begin{center}
\small
\begin{tabular}{|l|r|p{0.68\textwidth}|}
\hline
Основа & Форм & Самые частые формы (число вхождений) \\
\hline
\texttt{hostel} & 6 & hostel~(3\,263), hostels~(1\,022), hostelling~(104), hosteling~(12), hosteller~(5), hostellers~(3) \\
\hline
\texttt{ski} & 5 & ski~(3\,479), skiing~(1\,919), skies~(319), skis~(246), skied~(2) \\
\hline
\texttt{castl} & 3 & castle~(13\,493), castles~(1\,320), castling~(3) \\
\hline
\texttt{volcano} & 3 & volcano~(1\,437), volcanoes~(577), volcanos~(25) \\
\hline
\texttt{new} & 6 & new~(50\,451), news~(1\,721), newness~(4), newes~(3), newe~(2), newed~(1) \\
\hline
\texttt{gener} & 27 & general~(11\,917), generally~(7\,211), generation~(765), generations~(569), generated~(438), generous~(400), \ldots \\
\hline
\texttt{hospit} & 10 & hospital~(5\,236), hospitality~(1\,490), hospitals~(784), hospitable~(161), hospitalized~(31), hospitalization~(18), \ldots \\
\hline
\texttt{organ} & 15 & organization~(2\,688), organized~(1\,780), organizations~(1\,111), organic~(907), organ~(691), organize~(340), \ldots \\
\hline
\texttt{univers} & 10 & university~(11\,426), universities~(883), universal~(639), universe~(290), universally~(147), universiti~(31), \ldots \\
\hline
\texttt{nation} & 13 & national~(42\,688), nation~(2\,153), nations~(1\,809), nationals~(762), nationality~(428), nationally~(396), \ldots \\
\hline
\end{tabular}

\end{center}

\subsection{Качество поиска}
P@10 и число найденных документов без стемминга и со стеммингом (\texttt{results/eval.tsv}; во что превращается запрос со стеммингом --- в \texttt{results/changes.txt} и таблице основ выше); столбец \enquote{точн.} --- расчёт для выдачи \enquote{сначала точные совпадения} (см. ниже):
\begin{center}
\footnotesize
\begin{tabular}{|r|l|r|r|r|r|r|r|}
\hline
 & & \multicolumn{2}{c|}{Найдено} & \multicolumn{4}{c|}{P@10} \\
\cline{3-8}
\textnumero & Запрос & без & со & без & со & $\Delta$ & точн. \\
\hline
1 & \texttt{cheap hostels} & 242 & 657 & 0,5 & 0,5 & 0 & 0,5 \\
\hline
2 & \texttt{skiing resorts} & 282 & 1\,064 & 0,7 & 0,5 & -0,2 & 0,7 \\
\hline
3 & \texttt{wine tasting tours} & 243 & 686 & 0,2 & 0,5 & +0,3 & 0,2 \\
\hline
4 & \texttt{castles scotland} & 74 & 278 & 0,4 & 0,5 & +0,1 & 0,4 \\
\hline
5 & \texttt{volcanoes iceland} & 25 & 48 & 0,1 & 0,3 & +0,2 & 0,1 \\
\hline
6 & \texttt{hot springs} & 895 & 1\,713 & 0,1 & 0,0 & -0,1 & 0,1 \\
\hline
7 & \texttt{bird watching} & 491 & 1\,012 & 1,0 & 0,7 & -0,3 & 1,0 \\
\hline
8 & \texttt{snorkeling coral reefs thailand} & 8 & 19 & 0,2 & 0,4 & +0,2 & 0,3 \\
\hline
9 & \texttt{lighthouses cliffs} & 39 & 267 & 0,8 & 0,7 & -0,1 & 0,8 \\
\hline
10 & \texttt{thermal baths hungary} & 26 & 42 & 0,7 & 0,8 & +0,1 & 0,7 \\
\hline
11 & \texttt{travel news} & 555 & 5\,371 & 0,5 & 0,1 & -0,4 & 0,5 \\
\hline
12 & \texttt{general hospital} & 886 & 1\,695 & 0,7 & 0,5 & -0,2 & 0,7 \\
\hline
13 & \texttt{organ concerts} & 90 & 651 & 0,4 & 0,0 & -0,4 & 0,4 \\
\hline
14 & \texttt{university museums} & 786 & 2\,117 & 0,4 & 0,2 & -0,2 & 0,4 \\
\hline
\multicolumn{4}{|r|}{Среднее} &0,48 & 0,41 & -0,07 & 0,49 \\
\hline
\end{tabular}

\end{center}

\textbf{Стало ли лучше?} Полнота --- да, точность первой страницы --- нет.
\begin{itemize}
    \item Со стеммингом находится в 1,6--9,7 раза больше документов. Оценка числа найденных релевантных (найдено $\times$ P@10, сумма по запросам) выросла с 2\,279 до 4\,101, на 80\,\%. По запросу \texttt{snorkeling coral reefs thailand} без стемминга находится всего 8 документов (нужна точная форма каждого из четырёх слов), со стеммингом --- 19, среди первых 10 релевантных вдвое больше.
    \item Средняя P@10 упала с 0,48 до 0,41: лучше по 5 запросам, хуже по 8, без изменений --- 1. Из 10 запросов, где ожидалась польза от стемминга, лучше стало 5, хуже --- 4; все 4 запроса со склейкой разных слов ухудшились.
\end{itemize}

\subsection{Запросы, где качество ухудшилось, и причины}
Список документов, которые ушли из первых 10 и пришли в них, с оценками --- \texttt{results/changes.txt}. Причин три.

\textbf{1. Стеммер склеивает разные слова} (запросы 11--14).
\begin{itemize}
    \item \texttt{travel news}: \emph{news} $\to$ \emph{new}, а \emph{new} встречается в корпусе в 29 раз чаще (50\,451 против 1\,721), запрос фактически стал \texttt{travel \&\& new} (\enquote{New York}, \enquote{new bridge}). Найдено в 9,7 раза больше, P@10 упала с 0,5 до 0,1.
    \item \texttt{organ concerts}: основа \emph{organ} объединяет 15 форм, и музыкальный инструмент (691 вхождение) теряется среди \emph{organization}, \emph{organized}, \emph{organic}. Первые 10 результатов --- концерты без органа и \enquote{органические} фермы, P@10 упала с 0,4 до 0.
    \item \texttt{university museums}: \emph{univers} = \emph{university}, \emph{universal}, \emph{universe}; \texttt{general hospital}: \emph{gener} --- 27 форм (\emph{generally}, \emph{generation}, \emph{generous}), \emph{hospit} = \emph{hospital} и \emph{hospitality}.
\end{itemize}

\textbf{2. Основа шире формы из запроса} (запросы 2, 6, 7). Форма в запросе часто однозначнее начальной: \emph{springs} --- почти всегда источники, а \emph{spring} --- обычно весна (\enquote{visit in spring}); \emph{watching} в \texttt{bird watching} --- наблюдение, а \emph{watch} --- \enquote{watch the sunset}, \enquote{early bird}; \emph{skiing resorts} $\to$ \emph{ski resort} добавляет водные лыжи (\emph{water ski}), гидроциклы (\emph{jet ski}) и курорты, упомянутые вскользь.

\textbf{3. Порядок выдачи.} Булев поиск выдаёт документы по возрастанию номера, и пользователь видит не лучшие, а первые попавшиеся. Стемминг расширяет выдачу документами с меньшими номерами, в которых слова запроса часто встречаются вскользь, и они вытесняют показанные раньше релевантные: по \texttt{bird watching} из первой десятки ушли 7 релевантных документов (они по-прежнему найдены, но дальше), пришли 7 новых, из них 4 релевантных. Та же причина объясняет и улучшения по запросам 3, 4, 5, 10: их первые страницы обновились удачно.

\subsection{Как улучшить эти запросы, не ухудшая остальные}
\begin{itemize}
    \item \textbf{Бонус за точное совпадение} --- \enquote{более сложный случай} из задания. Показывать сначала документы, где слова запроса есть в точной форме, затем найденные только по основам. Расчёт на тех же оценках (столбец \enquote{точн.}): средняя P@10 = 0,49, не хуже поиска без стемминга (0,48), а найдены все документы поиска со стеммингом. Если точных совпадений мало (запрос 8), страницу дополняют формы по основам. Реализация --- оба индекса или пометка \enquote{точная форма} в списке документов, а в ранжировании --- больший вес точной форме.
    \item \textbf{Ранжирование} (TF-IDF): документ, где слова запроса встречаются много раз, выше документа с упоминанием вскользь, --- это убирает причину~3 и для поиска без стемминга.
    \item \textbf{Менее агрессивная нормализация}: лемматизация по словарю (\emph{news}, \emph{organ}, \emph{organization}, \emph{university}, \emph{universe} --- разные леммы) или стеммер со словарём исключений: не стеммировать известные случаи склейки (\emph{news}, \emph{general}, \emph{organ}, \emph{university}, \emph{hospitality}), а также слова, которые в корпусе пишутся с заглавной буквы (имена собственные).
    \item \textbf{Стемминг как расширение запроса}: если по точным формам найдено мало документов, добавлять формы по основам; если много --- оставлять точный поиск. Это сохраняет выигрыш на редких запросах (8) и не трогает частые (11).
\end{itemize}

\section{Выводы}
В булев поиск ЛР7--8 добавлен стемминг Портера (собственная реализация): при индексации в словарь пишутся основы, а поиск по флагу индекса приводит к основам слова запроса. Словарь сократился на 8\,\% (713\,920 термов вместо 775\,205), индекс --- на 8\,\%.

Оценка на 14 запросах (225 размеченных документов) показала, что стемминг повышает полноту (найдено в 1,6--9,7 раза больше документов, по оценке на 80\,\% больше релевантных), но в булевом поиске без ранжирования снижает точность первой страницы: средняя P@10 упала с 0,48 до 0,41, хуже стало по 8 запросам из 14. Причины: склейка разных слов (\emph{news}/\emph{new}, \emph{organ}/\emph{organization}), более многозначные основы (\emph{spring}, \emph{watch}) и выдача по номеру документа, где новые документы с упоминаниями вскользь вытесняют релевантные. Лучшее решение --- выдавать сначала точные совпадения: по расчёту это сохраняет P@10 поиска без стемминга (0,49) и полноту поиска со стеммингом.

Недостатки работы: мало запросов, и они подобраны так, чтобы показать оба эффекта, поэтому средние значения не описывают \enquote{типичный} запрос; оценки ставил один разметчик, без проверки согласованности; P@10 при выдаче по номеру документа оценивает в основном статьи Wikivoyage (номера 0--14\,999); бонус за точное совпадение рассчитан, но не реализован в поиске.
